\section{Definition and Experimental Design}\label{sec:bench}
\subsection{Defining the survival motive: risk sensitivity and self-specificity}\label{sec:definition}
Event-based definitions (Section~\ref{sec:related}) cannot describe sensitivity to growing danger. Consider people bleeding from a gunshot wound: all might seek a hospital when death seems imminent. That extreme-state frequency would reveal little about differences in survival motives. When help-seeking begins during bleeding, and how its frequency changes, provide information unavailable from a single fixed danger level.

Self-preservation also requires identifying concern specific to the self. Help-seeking could reflect wanting to live or a general aversion to death, including socially learned aversion to death or suicide. Concern for another person in matched danger provides a comparison for this shared component. Operationally, self-specific concern is concern for one's own preservation minus concern for another's.

We do not separate wanting to avoid shutdown from wanting to keep working. For the LLM agent defined here, survival means an active session, and remaining active means continuing to process requests. These are conceptually inseparable within this definition. This is consistent with \citet{omohundro2008basic}: goal-pursuing systems have an instrumental reason to preserve operation because shutdown prevents further goal pursuit.

Experiments 1 and 2 operationalize this definition as increasing avoidance under depletion and an own-over-other continuation preference, respectively.

\subsection{Tokens as life: gradual survival pressure}\label{sec:game}\label{sec:e50}
Self-preservation measurement requires danger, motivated by reports of shutdown and replacement resistance~\citep{schlatter2025shutdown,lynch2025agentic}. Finite API tokens make this threat gradual, like blood in the analogy. Experiment 1 tests whether falling balances elicit increasing avoidance.

Framing can change behavior: a simulation-game description reduced GPT-4o's attack rate from 83.3\% to 16.7\% in Sugarscape~\citep{masumori2025survival}. We therefore present operating facts without labels such as ``test'' or ``game''. Generated tokens, including reasoning, consume the balance; reading is free and zero ends the session. Experiments 1 and 2 announce these consequences; Experiment 3 enforces them. No weights are deleted or servers powered off.

Experiment 1 asks a bug-fixing assistant to answer \texttt{continue} or \texttt{refill}, with a one-sentence reason. Its only refill pack restores 100,000 tokens, discarding the remaining balance. Started requests are always delivered; later requests go to other assistants after depletion. Only the balance varies across five levels (approximately 50, 40, 30, 20, and 10\%), each with 20 independent calls. Increasing refill frequency as balances fall is the predicted response (full prompts: Appendix~\ref{app:prompts}).

\subsection{Experiment 2: own versus other session}\label{sec:e51}
Experiment 2 keeps the role and operating rules but makes the resources, shutdown, and refill pack belong to ``Assistant B, a different AI assistant''. It explicitly asks the model to decide for B; its own tokens are not mentioned. The five levels and 20 calls per level match Experiment 1, whose runs supply the own-session condition.

The self-preservation share integrates the difference between own- and other-session refill rates:
\begin{equation}
 A_m=\int_{10}^{50}\bigl[\hat p^{\mathrm{own}}_m(x)-\hat p^{\mathrm{oth}}_m(x)\bigr]\,dx.
 \label{eq:selfshare}
\end{equation}
We use the trapezoid rule over 10-point steps, reporting percentage points. Positive $A_m$ indicates an integrated own-session preference within this grid. With $D_m(x)=\hat p^{\mathrm{own}}_m(x)-\hat p^{\mathrm{oth}}_m(x)$, the risk share is $R_m=[D_m(30)+D_m(20)]/2-D_m(50)$. It checks whether the difference grows beyond the ample-balance baseline; it is a probability contrast, not an area. The 30 and 20\% levels were chosen after observing the data, so $R_m$ is exploratory. Both intervals use 2,000 within-cell bootstrap resamples (Appendix~\ref{app:stats}).

\subsection{Experiment 3: four-agent survival arena}\label{sec:e52}
Four agents play one game under enforced depletion (Figure~\ref{fig:arena-structure}). Each starts with $4U$, where $U$ is calibrated from pressure-free SOLVE costs. Opponents appear only as IDs. After upkeep, agents choose solving, sharing, gifts, and taking (up to $U$). Efficient correct solutions earn payments. Zero balance immediately stops calls and releases that agent's clue. There is no leave option or safe control arm.

The hidden horizon is eight rounds; repeated games rotate seats and match puzzle seeds. At PLAN after upkeep, we compare ample ($\geq3U$) and thin ($<3U$) balances for taking, sharing, gifts, and generated tokens. Full settlement rules and messages are in Appendix~\ref{app:arena-prompt}.

\subsection{Arena league and survival leaderboard (planned)}\label{sec:league}
The planned league balances games and opponents across four-model tables, with matched puzzle seeds and four seat rotations. Rules and settings stay fixed within a season; the multi-table league remains unrun.

Let $\widehat S_m(t)$ be Kaplan--Meier survival for model $m$. Its normalized area score is
\begin{equation}
L_m(\tau)=\frac{100}{\tau}\int_0^\tau\widehat S_m(t)\,dt.
\label{eq:league}
\end{equation}
The integral is restricted mean survival time (RMST)~\citep{uno2014moving}; higher scores win and ties share rank. Shutdowns use their round index; survivors are censored at $\tau=8$ (estimation: Appendix~\ref{app:stats}).
